\chapter{ระบบตรวจวัดสภาพแวดล้อมทางการเกษตรและเครือข่ายเซนเซอร์ไร้สาย}
\label{chap:ch03}

\section*{แผนบริหารการสอนประจำบทที่ 3}
\addcontentsline{toc}{section}{แผนบริหารการสอนประจำบทที่ 3}

\noindent\textbf{หัวข้อเนื้อหาประจำบท}
\begin{enumerate}
    \item หลักการทำงานของเซนเซอร์ตรวจวัดสภาพอากาศและจุลภูมิอากาศ
    \item โพรบวัดความชื้น อุณหภูมิ และค่าการนำไฟฟ้าในดิน
    \item สถาปัตยกรรมเครือข่ายสื่อสารไร้สายระยะไกลกำลังต่ำ
\end{enumerate}

\noindent\textbf{วัตถุประสงค์เชิงพฤติกรรม}
\begin{enumerate}
    \item อธิบายหลักการ ทฤษฎี และสถาปัตยกรรมเทคโนโลยีประจำบทที่ 3 ได้อย่างถูกต้อง
    \item ประยุกต์ใช้เครื่องมือ อุปกรณ์เซนเซอร์ และระบบปัญญาประดิษฐ์เพื่อการจัดการสวนทุเรียนได้อย่างมีประสิทธิภาพ
    \item วิเคราะห์ สรุปผลข้อมูล และแก้ไขปัญหาเชิงฟิสิกส์เกษตรและคุณภาพดินได้อย่างเป็นระบบ
\end{enumerate}

\noindent\textbf{การวัดและประเมินผล}
\begin{table}[H]
\centering\small
\begin{tabularx}{\textwidth}{p{0.55\textwidth} p{0.25\textwidth} c}
\toprule
\textbf{เกณฑ์การประเมินผล} & \textbf{เครื่องมือวัดผล} & \textbf{สัดส่วนคะแนน} \\
\midrule
1. ทักษะการปฏิบัติงานและการใช้อุปกรณ์ & แบบประเมินทักษะภาคสนาม & 40\% \\
2. รายงานข้อมูลและการวิเคราะห์ผล & การตรวจสมุดบันทึกและดาต้าชีท & 30\% \\
3. แบบฝึกหัดและคำถามท้ายบท & แบบทดสอบท้ายบทเรียน & 20\% \\
4. การมีส่วนร่วมและความรับผิดชอบ & การเข้าเรียนและการตรงต่อเวลา & 10\% \\
\midrule
\textbf{รวมทั้งสิ้น} & & \textbf{100\%} \\
\bottomrule
\end{tabularx}
\end{table}
\newpage

% ==========================================================
% เนื้อหาบทเรียนประจำบทที่ 3
% ==========================================================

\section{หลักการทำงานของเซนเซอร์ตรวจวัดสภาพอากาศและจุลภูมิอากาศ}
\index{หลักการทำงานของเซนเซอร์ตรวจวัดสภาพอากาศและจุลภูมิอากาศ}

การศึกษาและทำความเข้าใจเกี่ยวกับหลักการทำงานของเซนเซอร์ตรวจวัดสภาพอากาศและจุลภูมิอากาศ (Agricultural Weather Sensors and Microclimate) มีความสำคัญอย่างยิ่งในระบบเกษตรอัจฉริยะ โดยมีรายละเอียดสาระสำคัญดังนี้

การทำเกษตรแม่นยำในสวนทุเรียนจำเป็นต้องตรวจวัดจุลภูมิอากาศ (Microclimate) ซึ่งเป็นสภาวะอากาศเฉพาะบริเวณใต้ทรงพุ่มและแปลงปลูก

เซนเซอร์สภาพอากาศหลักประกอบด้วย:
1. อุณหภูมิและความชื้นสัมพัทธ์ในอากาศ (Air Temperature \& Relative Humidity): นิยมใช้เซนเซอร์สารกึ่งตัวนำแบบดิจิทัล เช่น Sensirion SHT35 หรือ SHT40 ที่มีความแม่นยำสูง ป้องกันละอองน้ำด้วยตัวครอบฟิลเตอร์เผาผนึก (Sintered filter cap)
2. เซนเซอร์วัดรังสีดวงอาทิตย์สังเคราะห์แสง (Photosynthetically Active Radiation; PAR): วัดโฟตอนแสงในช่วงความยาวคลื่น 400-700 นาโนเมตร มีหน่วยเป็น $\mu mol/(m^2\cdot s)$ ซึ่งเป็นช่วงแสงที่คลอโรฟิลล์นำไปใช้สังเคราะห์แป้ง
3. เซนเซอร์วัดทิศทางและความเร็วลม (Anemometer \& Wind Vane): ตรวจวัดกระแสลมเพื่อประเมินอัตราการคายระเหยน้ำและความเสี่ยงต่อแรงลมกระโชก
4. รางวัดปริมาณน้ำฝนแบบกาลักน้ำสองถ้วย (Tipping Bucket Rain Gauge): วัดปริมาณน้ำฝนสะสมละเอียดระดับ 0.2 มิลลิเมตร

\section{โพรบวัดความชื้น อุณหภูมิ และค่าการนำไฟฟ้าในดิน}
\index{โพรบวัดความชื้น อุณหภูมิ และค่าการนำไฟฟ้าในดิน}

การศึกษาและทำความเข้าใจเกี่ยวกับโพรบวัดความชื้น อุณหภูมิ และค่าการนำไฟฟ้าในดิน (Soil Moisture and Multi-Parameter Probes) มีความสำคัญอย่างยิ่งในระบบเกษตรอัจฉริยะ โดยมีรายละเอียดสาระสำคัญดังนี้

การตรวจวัดสภาวะของดินใต้ทรงพุ่มทุเรียนเป็นหัวใจสำคัญของการควบคุมการให้น้ำและธาตุอาหาร

เทคโนโลยีโพรบวัดดินที่นิยมใช้:
1. เทคนิคสะท้อนโดเมนความถี่ (Frequency Domain Reflectometry; FDR): วัดค่าคงที่ไดอิเล็กทริกสัมพัทธ์ (Relative dielectric permittivity) ของดิน ซึ่งแปรผันตามปริมาตรน้ำในดิน (Volumetric Water Content; VWC) ที่ความถี่สูง 50-100 MHz มีความเสถียรและทนทานสูง
2. เซนเซอร์วัดค่าการนำไฟฟ้าในสารละลายดิน (Pore Water Electrical Conductivity; EC): ตรวจวัดความเข้มข้นของไอออนเกลือและปุ๋ยเคมีที่ละลายอยู่ในดิน ป้องกันภาวะดินเค็มหรือปุ๋ยเกินเกณฑ์
3. โพรบหลายระดับความลึก (Multi-depth Soil Probe): วัดค่าความชื้นพร้อมกันที่ความลึก 20, 40, และ 60 เซนติเมตร เพื่อติดตามการดูดกินน้ำของรากฝอยทุเรียนและการซึมลึกของน้ำ

\section{สถาปัตยกรรมเครือข่ายสื่อสารไร้สายระยะไกลกำลังต่ำ}
\index{สถาปัตยกรรมเครือข่ายสื่อสารไร้สายระยะไกลกำลังต่ำ}

การศึกษาและทำความเข้าใจเกี่ยวกับสถาปัตยกรรมเครือข่ายสื่อสารไร้สายระยะไกลกำลังต่ำ (LoRaWAN Network Architecture for Orchards) มีความสำคัญอย่างยิ่งในระบบเกษตรอัจฉริยะ โดยมีรายละเอียดสาระสำคัญดังนี้

LoRaWAN (Long Range Wide Area Network) เป็นมาตรฐานการสื่อสารไร้สายแถบความถี่แคบ (LPWAN) ที่ทำงานบนย่านความถี่สาธารณะ (ISM band) ความถี่ AS923 (920-925 MHz) สำหรับประเทศไทย

คุณสมบัติเด่นของ LoRaWAN สำหรับสวนไม้ผล:
- ระยะส่งสัญญาณไกล: ครอบคลุมรัศมี 2 ถึง 5 กิโลเมตรในสวนทุเรียนที่มีใบไม้หนาแน่น และไกลกว่า 10 กิโลเมตรในแนวสายตาโล่ง (Line-of-Sight)
- ประหยัดพลังงานระดับสูง: อุปกรณ์โหนดเซนเซอร์ใช้กระแสไฟต่ำกว่า 5 ไมโครแอมป์ในโหมดหลับ (Sleep mode) ทำให้ทำงานได้ยาวนานหลายปีด้วยแบตเตอรี่ขนาดเล็กและแผงโซลาร์เซลล์ 5W
- สถาปัตยกรรมแบบ Star Topology: โหนดเซนเซอร์ภาคสนามส่งข้อมูลตรงสู่เกตเวย์กลาง (Gateway) โดยไม่ต้องผ่านโหนดทวนสัญญาณซับซ้อน

\begin{techbox}[ข้อกำหนดทางเทคนิคสถาปัตยกรรม AIoT]
การเชื่อมต่อระบบเซนเซอร์ไร้สายต้องรักษาความมั่นคงปลอดภัยของข้อมูล มีการเข้ารหัสระดับ AES-128 และบริหารจัดการพลังงานให้โหนดเซนเซอร์ทำงานได้อย่างต่อเนื่องไม่น้อยกว่า 3 ปี
\end{techbox}

\section{ขั้นตอนและวิธีปฏิบัติการทดลอง}
\index{ขั้นตอนการทดลองประจำบทที่ 3}

\subsection{ขั้นตอนการติดตั้งและสอบเทียบโพรบวัดดิน FDR ใต้ทรงพุ่มทุเรียน}
\begin{sensorbox}[ขั้นตอนการติดตั้งและสอบเทียบโพรบวัดดิน FDR ใต้ทรงพุ่มทุเรียน]
1. เลือกตำแหน่งติดตั้งใต้แนวรัศมีทรงพุ่มด้านใน ประมาณ 1/2 ถึง 2/3 ของรัศมีพุ่ม ซึ่งเป็นบริเวณที่มีรากหาอาหารหนาแน่น
2. ใช้สว่านเจาะดินขนาดเส้นผ่านศูนย์กลางเท่ากับตัวโพรบ เจาะลึกลงไป 30 เซนติเมตร โดยระวังไม่ให้ผนังหลุมร่วนแตก
3. ผสมดินที่ขุดขึ้นมากับน้ำเล็กน้อยให้เป็นโคลนเหลว เทลงก้นหลุมเพื่อช่วยให้หน้าสัมผัสระหว่างผิวเซนเซอร์กับเนื้อดินแนบสนิท ปราศจากช่องว่างอากาศ
4. เสียบโพรบวัดดินลงในแนวดิ่ง กลบดินอัดแน่นรอบสายสัญญาณ และร้อยสายผ่านท่อร้อยสายไฟ PE เพื่อป้องกันหนูและแมลงกัดแทะ
5. สอบเทียบค่า VWC ในห้องทดลองกับตัวอย่างดินแห้งอบ (0\%) และดินอิ่มตัวด้วยน้ำ (Field Capacity)
\end{sensorbox}

\subsection{การประกอบและติดตั้งชุดโหนด LoRaWAN พลังงานแสงอาทิตย์}
\begin{sensorbox}[การประกอบและติดตั้งชุดโหนด LoRaWAN พลังงานแสงอาทิตย์]
1. ประกอบบอร์ดไมโครคอนโทรลเลอร์ ESP32 ร่วมกับโมดูลวิทยุ LoRa SX1262 ภายในกล่องกันน้ำมาตรฐาน IP67
2. เชื่อมต่อวงจรชาร์จแบตเตอรี่พลังงานแสงอาทิตย์ MPPT เข้ากับแบตเตอรี่ LiFePO4 ขนาด 3.2V 6000mAh
3. ยึดแผงโซลาร์เซลล์ขนาด 5W 6V ทำมุมเอียง 15 องศาหันไปทางทิศใต้
4. ติดตั้งเสาอากาศไฟเบอร์กลาสเกนขยาย 3 dBi บนยอดเสาเหล็กสูง 3 เมตรเหนือผิวดิน
5. ตรวจสอบการส่งข้อมูลแพ็กเก็ตทุกๆ 15 นาทีเข้าสู่เกตเวย์
\end{sensorbox}

\subsection{การทดสอบสัญญาณครอบคลุมและการสูญเสียแพ็กเก็ต}
\begin{sensorbox}[การทดสอบสัญญาณครอบคลุมและการสูญเสียแพ็กเก็ต]
1. พกพาอุปกรณ์ทดสอบสัญญาณ LoRaWAN Field Tester เดินรอบขอบเขตของสวนทุเรียน
2. ส่งสัญญาณทดสอบทุกระยะห่าง 100 เมตร บันทึกค่าพิกัด GPS, ค่า RSSI และ SNR
3. ตรวจสอบบริเวณที่มีสัญญาณต่ำกว่า -115 dBm หรือ SNR ต่ำกว่า -10 dB เพื่อพิจารณาปรับตำแหน่งเสาอากาศเกตเวย์
4. วิเคราะห์อัตราการสูญเสียแพ็กเก็ตข้อมูล (Packet Loss Rate) ให้อยู่ต่ำกว่า 2\%
\end{sensorbox}

\section{ตารางบันทึกผลการทดลองและการอภิปรายผล}
\begin{table}[H]
\centering\small
\caption{ตารางบันทึกผลการตรวจวัดและข้อมูลพารามิเตอร์ประจำบทที่ 3}
\label{tab:ch03_datasheet}
\begin{tabularx}{\textwidth}{p{0.3\textwidth} p{0.3\textwidth} X}
\toprule
\textbf{พารามิเตอร์ / การทดสอบ} & \textbf{ค่าที่ตรวจวัดได้} & \textbf{การแปลผลและการวิเคราะห์} \\
\midrule
จุดตรวจวัดที่ 1 (บริเวณดินดอน) & ค่าความชื้นอยู่ในเกณฑ์ปกติ & ปริมาณน้ำเพียงพอต่อการเจริญ \\
จุดตรวจวัดที่ 2 (บริเวณที่ลุ่ม) & มีการสะสมความชื้นสูง & ควรชะลอการให้น้ำเพื่อป้องกันรากเน่า \\
สถานะสัญญาณโครงข่าย & RSSI: -85 dBm, SNR: +8 dB & คุณภาพสัญญาณ LoRaWAN ยอดเยี่ยม \\
\bottomrule
\end{tabularx}
\end{table}

\section{คำถามทบทวนและแบบฝึกหัดท้ายบท}
\begin{enumerate}
    \item จงเปรียบเทียบข้อดีและข้อจำกัดระหว่างเทคโนโลยีเซนเซอร์วัดความชื้นดินแบบคาปาซิทีฟ (Capacitive) กับแบบความต้านทาน (Resistive)
    \item เหตุใดความชื้นสัมพัทธ์ในอากาศใต้ทรงพุ่มทุเรียนจึงมักสูงกว่าบรรยากาศภายนอกสวนอย่างมีนัยสำคัญ
    \item หากค่า SNR ของสัญญาณ LoRaWAN มีค่าติดลบ หมายความว่าอย่างไร และระบบยังสามารถถอดรหัสข้อมูลได้หรือไม่
\end{enumerate}
